\input{estilo}
% Propuesta: la arquitectura interna en grilla. Arriba, el camino del publish;
% abajo, el del subscribe, con los dos frontends en la misma columna. La cola
% de retry debajo de la cola interna, la base debajo de los notifiers. Los
% pasos van numerados con una leyenda, y dos llaves marcan lo que el texto
% subraya: la misma pieza (cola + flota) que en SQS balancea, acá se usa para
% hacer fan-out. La frontera dice que es una hipótesis, como el texto.
\begin{document}
\begin{figura}\begin{tikzpicture}
  \newcommand\p[1]{\tikz[baseline=(x.base)]\node[paso] (x) {#1};}
  % la frontera
  \draw[draw=segundo, dashed, line width=0.9pt] (1.0,-5.75) rectangle (10.6,2.5);
  \node[text=segundo, anchor=north west] at (1.15,2.4) {SNS por dentro (hipótesis)};

  % --- el publish
  \draw[msg] (-0.4,0) -- node[cod, above] {publish} (1.3,0);
  \node[caja, minimum width=18mm, minimum height=12mm] (f) at (2.2,0) {frontend};
  \node[cola, minimum height=18mm] (q) at (5.4,0) {};
  \foreach \d in {0.3, 0.15} \node[caja, fill=white, minimum width=22mm, minimum height=14mm] at (8.8+\d,-\d) {};
  \node[caja, minimum width=22mm, minimum height=14mm] (w) at (8.8,0) {notifiers};
  \draw[msg] (f.east) -- node[paso, above=2pt] {2} (q.west);
  \draw[msg] (q.east) -- node[paso, above=2pt] {3} (w.west);
  \node[paso, fill=white] at (0.45,-0.35) {1};
  \foreach \i/\y in {1/1.3, 2/0, 3/-1.3} {
    \node[caja, minimum width=10mm, minimum height=10mm] (s\i) at (12.6,\y) {\var{S}\textsubscript{\i}};
    \draw[msg] (9.95,0) -- (s\i.west);
  }
  \node[paso, fill=white] at (11.05,-0.85) {5};
  \node[etiq, above] at (11.45,0) {sync};
  % las dos lecturas de la misma pieza
  \draw[llave] (4.4,1.05) -- (9.9,1.05) node[midway, above=6pt, etiq, text=suave] {cola + flota: balanceo de carga, como en SQS};
  \draw[llave] (10.8,1.95) -- (12.1,1.95) node[midway, above=6pt, etiq, text=suave] {fan-out};

  % --- los lentos
  \node[cola, minimum height=18mm] (r) at (5.4,-2.0) {retry};
  \draw[msg, draw=acento] (8.05,-0.75) -- node[paso, draw=acento, text=acento, fill=white, pos=0.45] {6} (6.35,-1.8);
  \draw[msg, draw=acento] (6.35,-2.2) -- (8.45,-0.75);

  % --- las suscripciones y el subscribe
  \pic[scale=1.6] at (8.8,-4.8) {base};
  \node[anchor=north] at (8.8,-5.1) {suscripciones};
  \draw[msg] (8.8,-0.75) -- node[paso, fill=white] {4} (8.8,-2.95);
  \draw[msg] (-0.4,-3.8) -- node[cod, above] {subscribe} (1.3,-3.8);
  \node[caja, minimum width=18mm, minimum height=12mm] (f2) at (2.2,-3.8) {frontend};
  \draw[msg] (f2.east) -- node[paso, above=2pt, pos=0.5] {\var{s}} (8.0,-3.8);

  % --- la leyenda
  \node[anchor=north west, align=left, font=\small] at (-0.4,-6.2) {%
    \p{1} \texttt{\small publish(tópico, m)}\\[2pt]
    \p{2} el frontend encola tópico y mensaje\\[2pt]
    \p{3} un notifier lo toma};
  \node[anchor=north west, align=left, font=\small] at (6.4,-6.2) {%
    \p{4} busca las suscripciones del tópico\\[2pt]
    \p{5} le manda \var{m} a cada suscriptor\\[2pt]
    \tikz[baseline=(x.base)]\node[paso, draw=acento, text=acento] (x) {6}; si uno no responde, a retry\\[2pt]
    \p{\var{s}} \texttt{\small subscribe}: agrega un registro};
\end{tikzpicture}\end{figura}
\end{document}
